\documentclass[12pt,reqno]{amsart}
\usepackage{math-review}
% Project-wide mathematical notation. No claims are certified here.

\title{Strong Compactness and Admissible Limits -- Lecture}
\author{Known-mathematics exposition benchmark}
\date{}
\begin{document}
\maketitle
\section{The goal and the reader's starting point}
Assume the reader can differentiate periodic functions, pass uniform limits
through integrals, extract a subsequence of bounded real numbers and diagonalize
a real symmetric matrix. We will build the stronger compactness interface needed here.

For an integer $d\geq1$, write $\rho$ for flat distance on
$\mathbb T^d=\mathbb R^d/\mathbb Z^d$ and set, for $0<\gamma<1$,
\[
 [v]_\gamma=\sup_{x\ne y}\frac{|v(x)-v(y)|}{\rho(x,y)^\gamma},\qquad
 \|u\|_{C^{2,\gamma}}
 =\max_{|\nu|\leq2}\|D^\nu u\|_\infty
   +\max_{|\nu|=2}[D^\nu u]_\gamma.
\]
Here $D^\nu$ denotes a periodic coordinate derivative of multi-index $\nu$.
Only the fixed flat structure is used below.

For matrices, $H\geq\delta I_r$ means $v^{\mathsf T}Hv\geq\delta|v|^2$
for every real vector $v$.
\begin{theorem}[Conditional sequential closedness]\label{thm:main}
Fix integers $d,r\geq1$ and $0<\beta<\alpha<1$.
On the flat torus $\mathbb T^d=\mathbb R^d/\mathbb Z^d$, let
\[
 A:C^{2,\beta}(\mathbb T^d)\longrightarrow
 C^0(\mathbb T^d;\operatorname{Sym}_r(\mathbb R))
\]
be continuous, where the target has the uniform operator norm, and set
\[
 \mathcal U=\{u\in C^{2,\beta}(\mathbb T^d):
       A(u)(x)>0\text{ for every }x\in\mathbb T^d\}.
\]
Let $\mathcal F:[0,1]\times\mathcal U\to C^0(\mathbb T^d)$ be continuous.
Suppose $t_j\in[0,1]$, $t_j\to t_\infty$, and
$u_j\in C^{2,\alpha}(\mathbb T^d)\cap\mathcal U$ satisfy, for constants
$C<\infty$ and $\delta>0$ independent of $j$,
\[
 \mathcal F(t_j,u_j)=0,\qquad
 \|u_j\|_{C^{2,\alpha}}\leq C,\qquad
 A(u_j)(x)\geq\delta I_r\quad\text{for every }j,x.
\]
Then a subsequence converges in $C^{2,\beta}$ to a function
$u_\infty\in C^{2,\alpha}\cap\mathcal U$, with
\[
 A(u_\infty)(x)\geq\delta I_r\quad\text{for every }x,
 \qquad \mathcal F(t_\infty,u_\infty)=0.
\]
The norms are the fixed periodic derivative/H\"older norms defined immediately
before this statement. This is a conditional sequential closedness assertion,
not an existence theorem or a smoothness assertion.
\end{theorem}
The route is to extract compatible derivatives, upgrade their convergence,
keep the limit admissible, and only then apply continuity of the equation.

\section{Building the compactness input}
The fixed $C^{2,\alpha}$ bound makes every derivative of order at most two
uniformly bounded and equicontinuous. Top derivatives have a common H\"older modulus;
lower derivatives have a common Lipschitz modulus from the next derivative.

Enumerate the rational points of the torus and the finitely many derivatives.
Repeated bounded-real subsequence extraction and a diagonal choice give convergence
at every listed point for every derivative.

For any one derivative sequence $v_j=D^\nu u_j$, a finite net upgrades
this to uniform convergence:
\[
 |v_j(x)-v_k(x)|
 \leq |v_j(x)-v_j(a)|+|v_j(a)-v_k(a)|+|v_k(a)-v_k(x)|.
\]
Choose a nearby net point $a$ so the first and third terms are uniformly small;
then convergence at the finitely many net points controls the middle term.
Apply this to all derivatives through order two along the same subsequence.

Denote the uniform derivative limits by $v_\nu$. They must still be shown to
be derivatives of the limiting function, not an unrelated collection. For $|\nu|\leq1$,
\[
 D^\nu u_j(x+he_i)-D^\nu u_j(x)
 =\int_0^h D^{\nu+e_i}u_j(x+se_i)\,ds.
\]
Uniform convergence passes through the integral. The resulting identity gives
$\partial_i v_\nu=v_{\nu+e_i}$. Thus $v_0=u_\infty\in C^2$, with exactly
these derivatives. Fixed-point difference quotients also retain their uniform
$\alpha$-H\"older bound, so $u_\infty\in C^{2,\alpha}$.

\section{Why the lower H\"older exponent gives strong convergence}
For an $\alpha$-H\"older function $w$, put $a=\|w\|_\infty$ and
$b=[w]_\alpha$. For points at distance $s>0$,
\[
 \frac{|w(x)-w(y)|}{s^\beta}
 \leq\min\{2a s^{-\beta},\,b s^{\alpha-\beta}\}.
\]
If $a,b>0$, split at $s_0=(2a/b)^{1/\alpha}$. The second bound applies
below $s_0$ and the first above it, giving
\[
 [w]_\beta\leq(2a)^{1-\beta/\alpha}b^{\beta/\alpha}.
\]
The cases $a=0$ or $b=0$ are immediate. Now set
$w=D^\nu(u_j-u_\infty)$ for each $|\nu|=2$. Its sup norm tends to zero,
while its $\alpha$ seminorm is at most the sum of two uniformly bounded seminorms.
All lower derivative sup norms already tend to zero. We have proved the exact input:
\begin{lemma}\label{lem:compactness}
On the flat torus $\mathbb T^d$ with fixed periodic derivative norms,
let $0<\beta<\alpha<1$. If a sequence of real functions satisfies
\[
 u_j\in C^{2,\alpha}(\mathbb T^d),\qquad
 \sup_j\|u_j\|_{C^{2,\alpha}}<\infty,
\]
then some subsequence converges strongly in $C^{2,\beta}(\mathbb T^d)$
to a function in $C^{2,\alpha}(\mathbb T^d)$.
The assertion does not claim convergence in the stronger $C^{2,\alpha}$ norm.
\end{lemma}
This is a proof of the compactness version used here, not merely a citation
of a general embedding theorem.

\section{Admissibility before continuity}
For the chosen subsequence, continuity of $A$ yields uniform convergence
$A(u_j)\to A(u_\infty)$. The relevant elementary matrix fact is the following.
\begin{lemma}\label{lem:margin}
Let $H_j,H\in C^0(\mathbb T^d;\operatorname{Sym}_r(\mathbb R))$,
let $\delta>0$ be fixed, and assume $H_j\to H$ uniformly in operator norm.
If $H_j(x)\geq\delta I_r$ for every $j,x$, then $H(x)\geq\delta I_r$
for every $x$. For positive definite $H_j(x)$ the same uniform lower bound
is equivalent to
\[
 \|H_j(x)^{-1}\|_{\mathrm{op}}\leq\delta^{-1}
 \quad\text{for every }j,x.
\]
\end{lemma}
For the limit assertion, fix $x$ and a unit vector $v$. Then
\[
 v^{\mathsf T}H(x)v=\lim_j v^{\mathsf T}H_j(x)v\geq\delta.
\]
This holds for every $x,v$, so $H\geq\delta I_r$. For the equivalent inverse
bound, orthogonal diagonalization gives
$\|H_j^{-1}\|_{\mathrm{op}}=1/\lambda_{\min}(H_j)$ at each point.

\begin{proof}[Proof of Theorem~\ref{thm:main}]
The compactness lemma gives a $C^{2,\beta}$ convergent subsequence with a
$C^{2,\alpha}$ limit. The preceding matrix argument proves
$A(u_\infty)\geq\delta I_r>0$, hence $u_\infty\in\mathcal U$.
Since $t_j\to t_\infty$ and $\mathcal F$ is continuous on this domain,
\[
 \mathcal F(t_\infty,u_\infty)
 =\lim_j\mathcal F(t_j,u_j)=0\quad\text{in }C^0.
\]
Each continuity hypothesis has now been used in its stated topology.
\end{proof}
No estimate was derived from a particular PDE: the norm bound and positive
margin were hypotheses. No regularity beyond $C^{2,\alpha}$ was established.

\section{Two computed limits of the method}
First take constant functions $u_j=t_j=1/j$, with $A(u)=uI_r$ and
$\mathcal F(t,u)=u-t$ on the open set of positive functions. These are
uniformly H\"older-bounded solutions, but $u_\infty=0$ is not admissible.
The missing hypothesis is one positive $\delta$ working for all $j$.

Second, set $u_n(x)=n^{-(2+\alpha)}\cos(2\pi n x_1)$. Since
$|\cos s-\cos t|\leq\min\{2,|s-t|\}$, the second derivatives satisfy
\[
 [\partial_1^2u_n]_\gamma\leq C_\gamma n^{\gamma-\alpha}
 \quad(0<\gamma<1),\qquad
 \|\partial_1^2u_n\|_\infty=(2\pi)^2n^{-\alpha}.
\]
Thus the $C^{2,\alpha}$ norms are bounded and the $C^{2,\beta}$ norms tend
to zero. Yet at $x_1=0,y_1=1/(2n)$ the $\alpha$ quotient is exactly
\[
 \frac{2(2\pi)^2n^{-\alpha}}{(1/(2n))^\alpha}
 =2(2\pi)^2 2^\alpha.
\]
So strong $C^{2,\alpha}$ convergence fails. The strict inequality
$\beta<\alpha$ has a mathematical purpose, not just a notational one.

\nocite{S-guide}
\bibliographystyle{plain}
\bibliography{../generated/references}
\end{document}
